% GENERATED FILE. Edit canonical lessons, then run npm run generate:books.
% canonical-source-hash: fnv1a64:c3f37547cf77280c
% canonical-lessons: SA-C173-ksipati, SA-C173-muncati, SA-C173-labhate, SA-C173-sevate, SA-C173-ramate

\chapter{Throws, Lets go, Gets, Serves, Enjoys}
\label{ch:sa-throws-lets-go-gets-serves-enjoys}
\hlchaptermodality{173}

\begin{chapteropening}
\textbf{By the end of this chapter:} \emph{I can say throws, lets go, gets, serves and enjoys.}

Complete the last lesson of chapter 173: I can say throws, lets go, gets, serves and enjoys.
\end{chapteropening}

\section[\texorpdfstring{\sk{क्षिपति} (kṣipati)}{क्षिपति (kṣipati)}]{\sk{क्षिपति} (kṣipati) — throws}

\label{lesson:SA-C173-ksipati}



\begin{quote}
Before the new one: say the Sanskrit for tourism, then the Sanskrit for a tourist.
\end{quote}

\subsection*{You'll want to know: \sk{क्षिपति}}
\textbf{\sk{क्षिपति}} — \emph{kṣipati} — \textquotedblleft{}throws\textquotedblright{}.

\begin{cousinweb}
\emph{kṣipati} comes from the root √kṣip \textquotedblleft{}throw, cast, strike\textquotedblright{} (CDIAL 3683), already in the Rigveda, and inherited from Indo-Iranian (EtymDB). \emph{nikṣipati} \textquotedblleft{}puts in\textquotedblright{} is built on the same root.
\end{cousinweb}

\begin{grammarlens}[title={What you've built}]
One more word for this chapter.
\end{grammarlens}

\subsection*{Your turn}
\begin{itemize}
\raggedright
  \item \emph{Say it:} \emph{kṣipati}
  \item \emph{Say it:} \emph{kṣipati}, once more
  \item \emph{Say it:} say \emph{kṣipati}
  \item \emph{Recall:} say the Sanskrit for tourism, then the Sanskrit for a tourist, then say \emph{kṣipati} again
\end{itemize}

\begin{tcolorbox}[breakable,colback=teal!4,colframe=teal!35!black,title={Before you move on}]
What does \sk{क्षिपति} mean? (\textquotedblleft{}Throws\textquotedblright{}.) Say it once more.
\end{tcolorbox}
\section[\texorpdfstring{\sk{मुञ्चति} (muñcati)}{मुञ्चति (muñcati)}]{\sk{मुञ्चति} (muñcati) — lets go, releases}

\label{lesson:SA-C173-muncati}



\begin{quote}
Before the new one: say the Sanskrit for a tourist, then the Sanskrit for throws.
\end{quote}

\subsection*{You'll want to know: \sk{मुञ्चति}}
\textbf{\sk{मुञ्चति}} — \emph{muñcati} — \textquotedblleft{}lets go, releases\textquotedblright{}.

\begin{cousinweb}
\emph{muñcati} comes from the root √muc \textquotedblleft{}release, let go, free\textquotedblright{} (CDIAL 10182), with a nasal \emph{ñ} added in the present, inherited from Indo-European (EtymDB). \emph{mokṣa} \textquotedblleft{}release, liberation\textquotedblright{} comes from the same root (MW).
\end{cousinweb}

\begin{grammarlens}[title={What you've built}]
One more word for this chapter.
\end{grammarlens}

\subsection*{Your turn}
\begin{itemize}
\raggedright
  \item \emph{Say it:} \emph{muñcati}
  \item \emph{Say it:} \emph{muñcati}, once more
  \item \emph{Say it:} say \emph{muñcati}
  \item \emph{Recall:} say the Sanskrit for a tourist, then the Sanskrit for throws, then say \emph{muñcati} again
\end{itemize}

\begin{tcolorbox}[breakable,colback=teal!4,colframe=teal!35!black,title={Before you move on}]
What does \sk{मुञ्चति} mean? (\textquotedblleft{}Lets go, releases\textquotedblright{}.) Say it once more.
\end{tcolorbox}
\section[\texorpdfstring{\sk{लभते} (labhate)}{लभते (labhate)}]{\sk{लभते} (labhate) — gets, receives}

\label{lesson:SA-C173-labhate}



\begin{quote}
Before the new one: say the Sanskrit for throws, then the Sanskrit for lets go.
\end{quote}

\subsection*{You'll want to know: \sk{लभते}}
\textbf{\sk{लभते}} — \emph{labhate} — \textquotedblleft{}gets, receives\textquotedblright{}.

\begin{cousinweb}
\emph{labhate} comes from the root √labh \textquotedblleft{}take, catch, obtain\textquotedblright{} (CDIAL 10948), inherited from Indo-European (EtymDB). MW compares it with Greek \emph{lambanō} \textquotedblleft{}take\textquotedblright{}, and \emph{lābha} \textquotedblleft{}profit\textquotedblright{} is its noun.
\end{cousinweb}

\begin{grammarlens}[title={What you've built}]
One more word for this chapter.
\end{grammarlens}

\subsection*{Your turn}
\begin{itemize}
\raggedright
  \item \emph{Say it:} \emph{labhate}
  \item \emph{Say it:} \emph{labhate}, once more
  \item \emph{Say it:} say \emph{labhate}
  \item \emph{Recall:} say the Sanskrit for throws, then the Sanskrit for lets go, then say \emph{labhate} again
\end{itemize}

\begin{tcolorbox}[breakable,colback=teal!4,colframe=teal!35!black,title={Before you move on}]
What does \sk{लभते} mean? (\textquotedblleft{}Gets, receives\textquotedblright{}.) Say it once more.
\end{tcolorbox}
\section[\texorpdfstring{\sk{सेवते} (sevate)}{सेवते (sevate)}]{\sk{सेवते} (sevate) — serves, looks after}

\label{lesson:SA-C173-sevate}



\begin{quote}
Before the new one: say the Sanskrit for lets go, then the Sanskrit for gets.
\end{quote}

\subsection*{You'll want to know: \sk{सेवते}}
\textbf{\sk{सेवते}} — \emph{sevate} — \textquotedblleft{}serves, looks after\textquotedblright{}.

\begin{cousinweb}
\emph{sevate} comes from the root √sev \textquotedblleft{}associate with, attend on\textquotedblright{} (CDIAL 13593), already in the Rigveda. \emph{sevā} \textquotedblleft{}service\textquotedblright{} and \emph{sevaka} \textquotedblleft{}a servant\textquotedblright{} come from the same root.
\end{cousinweb}

\begin{grammarlens}[title={What you've built}]
One more word for this chapter.
\end{grammarlens}

\subsection*{Your turn}
\begin{itemize}
\raggedright
  \item \emph{Say it:} \emph{sevate}
  \item \emph{Say it:} \emph{sevate}, once more
  \item \emph{Say it:} say \emph{sevate}
  \item \emph{Recall:} say the Sanskrit for lets go, then the Sanskrit for gets, then say \emph{sevate} again
\end{itemize}

\begin{tcolorbox}[breakable,colback=teal!4,colframe=teal!35!black,title={Before you move on}]
What does \sk{सेवते} mean? (\textquotedblleft{}Serves, looks after\textquotedblright{}.) Say it once more.
\end{tcolorbox}
\section[\texorpdfstring{\sk{रमते} (ramate)}{रमते (ramate)}]{\sk{रमते} (ramate) — enjoys, delights in}

\label{lesson:SA-C173-ramate}



\begin{quote}
Before the new one: say the Sanskrit for gets, then the Sanskrit for serves.
\end{quote}

\subsection*{You'll want to know: \sk{रमते}}
\textbf{\sk{रमते}} — \emph{ramate} — \textquotedblleft{}enjoys, delights in\textquotedblright{}.

\begin{cousinweb}
\emph{ramate} comes from the root √ram \textquotedblleft{}stop, rest, take pleasure in\textquotedblright{} (CDIAL 10628), inherited from Indo-European (EtymDB); MW compares Greek \emph{eramai} \textquotedblleft{}love\textquotedblright{}. It lives on as Gujarati \emph{ramvũ} \textquotedblleft{}to enjoy oneself, saunter about\textquotedblright{} (CDIAL 10628).
\end{cousinweb}

\begin{grammarlens}[title={What you've built}]
One more word for this chapter.
\end{grammarlens}

\subsection*{Your turn}
\begin{itemize}
\raggedright
  \item \emph{Say it:} \emph{ramate}
  \item \emph{Say it:} \emph{ramate}, once more
  \item \emph{Say it:} say \emph{ramate}
  \item \emph{Recall:} say the Sanskrit for gets, then the Sanskrit for serves, then say \emph{ramate} again
\end{itemize}

\begin{tcolorbox}[breakable,colback=teal!4,colframe=teal!35!black,title={Before you move on}]
What does \sk{रमते} mean? (\textquotedblleft{}Enjoys, delights in\textquotedblright{}.) Say it once more.
\end{tcolorbox}
